\documentclass[multi=tikzpicture,border=9pt]{standalone}
\usepackage[UTF8]{ctex}
\usepackage{amsmath,amssymb}
\usepackage{tikz}
\usetikzlibrary{arrows.meta,calc,patterns,decorations.pathmorphing}
\definecolor{main}{HTML}{17365D}
\definecolor{accent}{HTML}{AC3931}
\tikzset{every picture/.style={font=\small,line width=.75pt},
>={Stealth[length=2mm]},arr/.style={->},box/.style={draw=main,rounded corners=2pt,
fill=main!4,text width=2.7cm,minimum height=.9cm,align=center},
note/.style={align=center,text width=8cm}}
\begin{document}
\begin{tikzpicture}
\draw[fill=main!10] (-3,1.1) rectangle (2,2.2);\node at (-.5,1.65){巨石 $M$};\draw[arr,accent] (2,1.65)--(3.2,1.65) node[right]{$F$};\draw[thick] (-3.5,0)--(3.5,0);\draw (-2.4,.55) circle(.55);\draw (-1.2,.55) circle(.55);\draw (0,.55) circle(.55);\draw (1.2,.55) circle(.55);\node at (-.8,-.55){10根圆木（数量示意）};\begin{scope}[xshift=6cm,yshift=1.3cm]\draw (0,0) circle(1.4);\draw[thick] (-2,1.3)--(2,1.3);\draw[thick] (-2,-1.3)--(2,-1.3);\fill (-.52,1.3) circle(2pt) node[above]{$A$};\fill (.52,-1.3) circle(2pt) node[below]{$B$};\draw[dashed] (-.52,1.3)--(.52,-1.3);\fill (0,0) circle(1pt);\node[right] at (0,0){$O$};\draw[dashed] (0,-1.3)--(0,1.3);\draw (0,.6) arc(90:112:.6);\node at (-.25,.75){$\theta$};\end{scope}\node[note] at (2,-1.7){压缩后竖直间距为 $2r(1-\eta)$。图中偏移角夸大以便辨认。};
\end{tikzpicture}
\begin{tikzpicture}
\draw (0,0) circle(1.4);\draw[thick] (-2,1.3)--(2,1.3);\draw[thick] (-2,-1.3)--(2,-1.3);\fill (-.52,1.3) circle(2pt) node[above]{$A$};\fill (.52,-1.3) circle(2pt) node[below]{$B$};\draw[dashed] (-.52,1.3)--(.52,-1.3);\fill (0,0) circle(1pt);\node[right] at (0,0){$O$};\draw[arr,accent] (-.52,1.3)--(-.52,.45) node[left]{$N$};\draw[arr,accent] (.52,-1.3)--(.52,-.45) node[right]{$N$};\draw[arr,main] (-.52,1.3)--(.7,1.3) node[right]{$f$};\draw[arr,main] (.52,-1.3)--(-.7,-1.3) node[left]{$f$};\draw[dashed] (0,-1.3)--(0,1.3);\node[note] at (0,-2.5){压力偏移力矩 $2Nr\sin\theta$；切向力矩 $2fr\cos\theta$，两者相反。};
\end{tikzpicture}
\begin{tikzpicture}
\node[box] (b0) at (0,0) {单根压力 $N$};
\node[box] (b1) at (5,0) {$2Nr\sin\theta$\\$=2fr\cos\theta$};
\node[box] (b2) at (5,-2) {$f/N=\tan\theta$};
\node[box] (b3) at (0,-2) {$F=\sum f_i$\\$=Mg\tan\theta$};
\draw[arr] (b0)--(b1);\draw[arr] (b1)--(b2);\draw[arr] (b2)--(b3);\node[note] at (2.5,-3.35) {压力与切向力各自合力为零，但两对力矩要相互抵消。};
\end{tikzpicture}
\begin{tikzpicture}
\draw (0,0) circle(1.4);\draw[thick] (-2,1.3)--(2,1.3);\draw[thick] (-2,-1.3)--(2,-1.3);\fill (-.52,1.3) circle(2pt) node[above]{$A$};\fill (.52,-1.3) circle(2pt) node[below]{$B$};\draw[dashed] (-.52,1.3)--(.52,-1.3);\fill (0,0) circle(1pt);\node[right] at (0,0){$O$};\draw[arr,main] (0,0)--(1.1,0) node[right]{$V$};\draw[arr,accent] (-.52,1.7)--(1.2,1.7) node[right]{$v_{\rm stone}=2V$};\draw[arr] (1.1,.6) arc(30:-65:1.1);\node[right] at (1.35,-.45){$\omega$};\node[note] at (0,-2.5){下接触：$V-\omega r\cos\theta=0$；上接触：$V+\omega r\cos\theta=2V$。};
\end{tikzpicture}
\begin{tikzpicture}
\node[box] (b0) at (0,0) {$\eta=0.002$};
\node[box] (b1) at (5,0) {$\cos\theta=1-\eta$};
\node[box] (b2) at (5,-2) {$F=Mg\tan\theta$};
\node[box] (b3) at (0,-2) {两处均需\\$\mu_s\geq\tan\theta$};
\draw[arr] (b0)--(b1);\draw[arr] (b1)--(b2);\draw[arr] (b2)--(b3);\node[note] at (2.5,-3.35) {结果依赖题设压力合力偏移模型。无滑动不等于切向力为零；形状恢复不等于能量无损。};
\end{tikzpicture}
\end{document}
